\section{Exact power pooling and integral algebra}
\subsection{The inherited two row certificate}
For $|A|>1$ let $(\chi_A(k),\psi_A(k))$, $k\in\Z$, be the canonical
signed Pell pair. Its recurrence matrix is
$\left(\begin{smallmatrix}A&A^2-1\\1&A\end{smallmatrix}\right)$,
with initial pair $(1,0)$. Its determinant is one, so all signed indices
are defined. Given a nonzero base $n$, put
\[
 T_n=2An-n^2-1.
\]
Then $\gcd(n,T_n)=1$, and
\begin{equation}\label{eq:pooled-congruence}
 \chi_A(k)+(n-A)\psi_A(k)\equiv n^k\pmod{T_n}
 \qquad(k\in\Z).
\end{equation}
Indeed $(n-A)^2\equiv A^2-1\pmod{T_n}$; evaluating the fundamental
unit at $n-A$ gives $n$. Positive powers follow by multiplication,
and negative powers by invertibility. The Pell functions are proof
notation, not primitives in the final polynomial.

For $r$ requests $(k_i,m_i)$ at this base, retain the existing coordinates
$X_i=\chi_A(k_i)$, $Y_i=\psi_A(k_i)$ and $Z_i=AY_i$.
Introduce only two new integer coordinates $K,L$, and impose
\begin{align}
 K-2\sum_iY_i(m_i-X_i)-(2A-n)L&=0,\label{eq:pool-row1}\\
 \sum_i\bigl((m_i-X_i+Z_i)^2-Y_i^2\bigr)-nK+L&=0.
 \label{eq:pool-row2}
\end{align}
Every term has degree at most two. Neither $T_n$ nor any displayed
summation abbreviation introduces another independent coordinate.

\begin{lemma}\label{lem:pooled-identity}
The two rows imply
\[
 \sum_i\bigl(m_i-X_i+(A-n)Y_i\bigr)^2
 =T_n\left(L-\sum_iY_i^2\right).
\]
Conversely, when each error on the left is divisible by nonzero $T_n$,
there are integers $K,L$ satisfying both rows.
\end{lemma}
\begin{proof}
Set $D_i=m_i-X_i+AY_i$ and $H=\sum_iY_i(m_i-X_i)$.
Eliminating $K$ gives
$\sum_i(D_i^2-Y_i^2)-2nH=T_nL$.
Expanding $\sum_i(D_i-nY_i)^2$ subtracts exactly
$T_n\sum_iY_i^2$, proving the identity. For the converse choose
\[
 L=\sum_iY_i^2+T_n^{-1}\sum_i(D_i-nY_i)^2,
 \qquad K=2H+(2A-n)L.
\]
The divisibility assumption makes $L$ integral and direct substitution
proves both rows. No sign of $K$ or $L$ is required.
\end{proof}

\begin{theorem}[Simultaneous exactness from a square frame]
\label{thm:pooled-exactness}
Suppose $W\ge258$, $1\le r\le15$, $n\ne0$, $|A|>W^W$, and
$n^2,k_i^2,m_i^2<W$. If the pairs are the specified canonical signed
Pell pairs, the two quadratic rows force $m_i=n^{k_i}$ for every $i$.
For $|n|\ge2$ they also force $k_i\ge0$; for $n=\pm1$ all signed
indices are allowed. Conversely all these exact powers admit $K,L$.
\end{theorem}
\begin{proof}
Put $b=\lfloor\sqrt{W-1}\rfloor\ge16$. Every absolute base, index
and proposed output is at most $b$. Multiply the aggregate congruence
by $n^{2b}$, using \eqref{eq:pooled-congruence}, to get
\[
 T_n\mid E,\qquad
 E=\sum_i(m_in^b-n^{b+k_i})^2\ge0.
\]
All exponents are nonnegative. Since $b\ge2$,
$|m_in^b-n^{b+k_i}|\le b^{b+1}+b^{2b}\le2b^{2b}$.
Consequently
\[
 E\le4rb^{4b}<W^{2b+1}\le W^W,
\]
because $4r\le60<W$, $b^2<W$ and $2b+1\le b^2+1\le W$.
On the other hand
$|T_n|\ge2|A|-W>|A|>W^W$. Divisibility forces $E=0$,
so every summand vanishes. A nonnegative $k_i$ gives the asserted power
by cancelling $n^b$. If $k_i=-j<0$, cancellation gives $m_in^j=1$;
this excludes nonunit bases and gives the stated power at the two units.
For the converse each exact error is divisible by $T_n$, so apply
Lemma~\ref{lem:pooled-identity}.
\end{proof}
The absolute-value argument includes negative radices and the unit case;
assuming a positive radix before decoding its sign would not suffice.
The actual square-frame ledger, canonical-pair classification and
request ports are bound to the literal source by the finite checker.
A different affine frame that only gives $|a|<W$ cannot silently be
substituted for this square-frame hypothesis.

\subsection{Exact integral projections along quadratic radicals}
Write each residual as $Q_i=q_i+\ell_i+c_i$, with homogeneous quadratic
part $q_i$, linear part $\ell_i$, and integer Hessian $H_i$.
For an integer vector $v$ with $H_iv=0$ for every $i$, exact expansion gives
\begin{equation}\label{eq:radical-slope}
 Q_i(y+sv)=Q_i(y)+s\ell_i(v).
\end{equation}
The mixed term is $sy^TH_iv=0$ and the pure term vanishes since
$2q_i(v)=v^TH_iv=0$.

\begin{theorem}\label{thm:radical-projection}
Suppose $v_p=1$ and one slope $a_r=\ell_r(v)$ is $\pm1$.
With $y_p=0$, let
\[
 R_i(y)=Q_i(y)-a_ia_rQ_r(y),\quad i\ne r,
 \qquad a_i=\ell_i(v).
\]
These integer quadratics have a common zero set in bijection with that
of all $Q_i$, with one fewer variable and residual. The lift is
$x=y-a_rQ_r(y)v$, and the inverse is $y=x-x_pv$.
\end{theorem}
\begin{proof}
Put $s=-a_rQ_r(y)$. Equation~\eqref{eq:radical-slope} gives
$Q_r(y+sv)=0$ and $Q_i(y+sv)=R_i(y)$ for $i\ne r$.
Conversely an old zero gives $y_p=0$ and
$Q_r(y)=-x_pa_r$, so the displayed lift restores precisely $x_p$.
All other $R_i(y)$ vanish by the same formula. The two maps are inverse
and integral; the new rows are constant combinations of restricted
quadratics. No division condition remains.
\end{proof}
A primitive direction alone would not justify division by a nonunit
slope. The actual certificates require the two literal unit conditions.
These unrestricted-integer coordinate changes need not preserve a
positive domain.

\begin{lemma}[Affine consequence elimination]\label{lem:affine-elimination}
Suppose a nonzero constant integer combination $L=\sum_i\lambda_iQ_i$
has degree at most one and its primitive affine normalization has
coefficient $\pm1$ on $x_p$. Solve it as $x_p=S(\widehat{x_p})$.
Substitute in every row and omit any row $r$ with $\lambda_r\ne0$.
The resulting quadratic system has an integer zero set in bijection
with the original system.
\end{lemma}
\begin{proof}
Every original zero satisfies $L=0$, hence $x_p=S$ and every retained
substituted row. Conversely, the substitution makes $L$ identically
zero. At a zero of the retained rows, $\lambda_rQ_r=0$; the integers
have no zero divisors, so $Q_r=0$. Insertion of $S$ and deletion of
$x_p$ are inverse. A unit $\lambda_r$ is not needed; a unit solved
variable coefficient after primitive normalization is essential.
\end{proof}

An elementary integer row operation $Q_i\leftarrow Q_i+kQ_j$,
$i\ne j$, has inverse with $-k$. A triangular integral affine coordinate
change has its explicit integral inverse. A monic quadratic definition
can be eliminated only when its complete substitution retains the claimed
degree. Each released certificate records these conditions and the exact
index compaction. A sum of squares of the final rows is generally a
\emph{different} polynomial, with the same zero-existence statement.

